% Códigos e saídas: fonte um pouco menor e quebra automática das linhas
% que excederem a largura do texto (pacote fvextra).
% Se o fvextra não estiver instalado, o livro ainda compila, mas sem a
% quebra automática (linhas longas podem ultrapassar a margem).
\IfFileExists{fvextra.sty}{%
  \usepackage{fvextra}
  \DefineVerbatimEnvironment{Highlighting}{Verbatim}{commandchars=\\\{\},fontsize=\small,breaklines,breakanywhere,breaknonspaceingroup}
  \RecustomVerbatimEnvironment{verbatim}{Verbatim}{fontsize=\small,breaklines,breakanywhere}
}{%
  \PackageWarningNoLine{MCBS}{fvextra.sty nao encontrado: codigos sem quebra automatica de linha}
  \DefineVerbatimEnvironment{Highlighting}{Verbatim}{commandchars=\\\{\},fontsize=\small}
  \RecustomVerbatimEnvironment{verbatim}{Verbatim}{fontsize=\small}
}
% Permite mais espaçamento entre palavras antes de deixar uma linha
% ultrapassar a margem (útil em parágrafos com trechos de código).
\setlength{\emergencystretch}{3em}
% Títulos de nível 4 e 5 (\paragraph e \subparagraph) em linha própria;
% no padrão da classe book, eles ficam na mesma linha do texto seguinte e,
% quando seguidos de código, aparecem com espaços esticados.
\titleformat{\paragraph}[hang]{\normalfont\normalsize\bfseries}{\theparagraph}{1em}{}
\titlespacing*{\paragraph}{0pt}{3.25ex plus 1ex minus .2ex}{1ex plus .2ex}
\titleformat{\subparagraph}[hang]{\normalfont\normalsize\bfseries}{\thesubparagraph}{1em}{}
\titlespacing*{\subparagraph}{0pt}{3.25ex plus 1ex minus .2ex}{1ex plus .2ex}
